\documentclass[10pt,a4paper]{amsart}
\usepackage[margin=19mm]{geometry}
\usepackage{amsmath,amssymb,mathtools}
\usepackage[scr=rsfs,cal=euler]{mathalfa}
\usepackage{xfrac}
\usepackage[unicode,hidelinks,bookmarks=false]{hyperref}
\newcommand{\ie}{i.e.\ }
\newcommand{\cf}{cf.\ }
\newcommand{\resp}{respectively\ }
\newcommand{\Subsection}{\subsection}
\providecommand{\nobreakdash}{\nobreak\hbox{-}}
\setlength{\emergencystretch}{2em}
\multlinegap0pt
\usepackage{lmodern}
\usepackage{mathtools}
\usepackage{enumitem}
\usepackage{tikz-cd}
\newtheorem{theorem}{Theorem}[section]
\newtheorem{proposition}[theorem]{Proposition}
\newtheorem{lemma}[theorem]{Lemma}
\newtheorem{corollary}[theorem]{Corollary}
\newtheorem{question}[theorem]{Question}
\newtheorem{conjecture}[theorem]{Conjecture}
\newtheorem{definition}[theorem]{Definition}
\newtheorem{example}[theorem]{Example}
\newtheorem{notation}[theorem]{Notation}
\newtheorem{remark}[theorem]{Remark}
\newcommand{\coker}{\operatorname{coker}}
\newcommand{\A}{\mathcal A}
\newcommand{\W}{\mathcal W}
\newcommand{\M}{\mathcal M}
\newcommand{\B}{\mathcal B}
\newcommand{\E}{\mathcal E}
\newcommand{\T}{\mathsf T}
\newcommand{\F}{\mathsf F}
\newcommand{\tors}{\operatorname{tors}}
\newcommand{\torf}{\operatorname{torf}}
\newcommand{\filt}{\operatorname{filt}}
\newcommand{\gen}{\operatorname{gen}}
\newcommand{\cogen}{\operatorname{cogen}}
\newcommand{\add}{\operatorname{add}}
\newcommand{\simp}{\operatorname{simp}}
\newcommand{\Hom}{\operatorname{Hom}}
\newcommand{\Ext}{\operatorname{Ext}}
\newcommand{\Up}{\operatorname{Up}}
\newcommand{\Low}{\operatorname{Low}}
\newcommand{\id}{\operatorname{id}}
\newcommand{\elll}{\ell}
\newcommand{\perpzero}{{\perp_0}}
\newcommand{\lperpzero}{{{}^{\perp_0}}}
\newcommand{\fCogen}[1]{\mathbf{F}(#1)}
\newcommand{\tGen}[1]{\mathbf{T}(#1)}
\newcommand{\tpair}[2]{(\mathbf{#1}, \mathbf{#2})  }
\begin{document}
\title[Wide subcategories and brick-finiteness for length categorie]{Wide subcategories and brick-finiteness for length categories}
\author{Francesco Sentieri}
\date{}
\maketitle
\noindent\emph{Source and licence.} Wide subcategories and brick-finiteness for length categories. Version arXiv:2607.15991v1 (CC BY 4.0). This material is distributed under the Creative Commons Attribution 4.0 International licence, http://creativecommons.org/licenses/by/4.0/, as recorded in the arXiv OAI licence metadata for this exact version and shown on the abstract page https://arxiv.org/abs/2607.15991v1. Full rights notice: licenses/arxiv-2607.15991-CC-BY-4.0.txt.\par\smallskip
\noindent\textit{Editorial scope.} Counted unit: the original \texttt{theorem} environment \texttt{thm:full-mutation-labels} at source lines 615--627 together with its complete proof at lines 629--660; the delivered document keeps source lines 210--661 so that every referenced label and every citation resolves inside it. Native editable source is the paper's own concatenated TeX; the AMS wrapper is editorial formatting only. No publisher class, upstream script or unrelated artwork is redistributed.
\begin{proposition}[{\cite[Lemma 4.4]{EnomotoMonobrick}}]\label{prop:maximal-f-simple}
An \(\mathbf f\)-simple object \(B\) is torsionfree, almost torsion if and only if, for every \(\mathbf f\)-simple object \(B'\), we have
\[
\Hom_\A(B,B')\neq 0
\quad\Longleftrightarrow\quad
B'\simeq B.
\]
In particular, the torsionfree, almost torsion objects for a torsion pair form a semibrick. Dually, the torsion, almost torsionfree objects form a semibrick.
\end{proposition}

We will use the following Ext-orthogonality several times. 

\begin{lemma}\label{lem:ext-orthogonality}
Let \((\mathbf t,\mathbf f)\) be a torsion pair, let \(T\) be torsion, almost torsionfree, and let \(F\) be torsionfree, almost torsion for this torsion pair. Then
\[
\Ext^1_\A(T,F)=0.
\]
\end{lemma}

\begin{proof}
We give a short proof for the convenience of the reader. Let 
\[
0 \to F \to M \to T \to 0
\]
be a short exact sequence. Then, if the sequence were not split, as $ F $ is torsionfree, almost torsion, $ M $ would have to be in $ \mathbf{t} $. But then, as $ T $ is torsion, almost torsionfree, we would have that $ F \in \mathbf{t} $. This would imply $ F = 0 $, which is a contradiction.
\end{proof}


\subsection{Wide subcategories and compactness}

\begin{definition}
A subcategory \(\W\subseteq\A\) is wide if it is closed under kernels, cokernels and extensions.

A torsion pair \((\mathbf t,\mathbf f)\) is widely cogenerated if \(\mathbf f=\F(\W)\) for some wide subcategory \(\W\). Dually, it is widely generated if \(\mathbf t=\T(\W)\) for some wide subcategory \(\W\).
\end{definition}

\begin{theorem}[{\cite[Section 1.2]{RingelSpecies}}]\label{thm:ringel-semibricks}
There is a bijection between semibricks in \(\A\) and wide subcategories of \(\A\). This bijection associates with a semibrick \(\mathcal S\) the extension closure \(\filt(\mathcal S)\). Its inverse associates with every wide subcategory the collection of its simple objects.
\end{theorem}

\begin{proposition}[{\cite[Proposition 5.2]{EnomotoMonobrick}}]\label{prop:wide-cogenerated-criterion}
Let \((\mathbf t,\mathbf f)\) be a torsion pair. The following statements are equivalent.
\begin{enumerate}[label=\textup{(\arabic*)},leftmargin=2.5em]
\item \((\mathbf t,\mathbf f)\) is widely cogenerated.
\item Every \(\mathbf f\)-simple object is a subobject of some torsionfree, almost torsion object for \((\mathbf t,\mathbf f)\).
\item \((\mathbf t,\mathbf f)\) is widely cogenerated by the extension closure of its torsionfree, almost torsion objects.
\end{enumerate}
Dually, \((\mathbf t,\mathbf f)\) is widely generated if and only if every \(\mathbf t\)-simple object is a quotient of some torsion, almost torsionfree object.
\end{proposition}

\begin{proposition}[{\cite[Corollary 5.6]{RingelBrickChain}}]\label{prop:finitely-cogenerated}
Let \((\mathbf t,\mathbf f)\) be a torsion pair. Then \(\mathbf f\) is finitely cogenerated, that is \(\mathbf f=\F(M)\) for some object \(M\), if and only if \((\mathbf t,\mathbf f)\) is widely cogenerated and has only finitely many torsionfree, almost torsion objects.

Dually, \(\mathbf t\) is finitely generated if and only if \((\mathbf t,\mathbf f)\) is widely generated and has only finitely many torsion, almost torsionfree objects.
\end{proposition}

\begin{definition}
We say that \(\A\) is widely determined, or WD, if every torsion class in \(\A\) is widely generated. We say that \(\A\) is widely co-determined, or WCD, if every torsionfree class in \(\A\) is widely cogenerated.
\end{definition}

\begin{corollary}\label{cor:basic-wide-consequences}
Assume \(\A\) is WD and WCD.
\begin{enumerate}[label=\textup{(\arabic*)},leftmargin=2.5em]
\item Every non-zero torsion class has a torsion, almost torsionfree object.
\item If a torsion class is not compact, then it has infinitely non-isomorphic torsion, almost torsionfree objects.
\item If a torsion class is not co-compact, then it has infinitely many non-isomorphic torsionfree, almost torsion objects.
\end{enumerate}
\end{corollary}

\begin{proof}
For (1), let \(\mathbf t\neq 0\) and write \(\mathbf t=\T(\W)\) with \(\W\) wide. The wide subcategory \(\W\) is non-zero, hence has a simple object. By Theorem \ref{thm:ringel-semibricks} and the dual of Proposition \ref{prop:wide-cogenerated-criterion}, this gives a torsion, almost torsionfree object.

For (2), if there were only finitely many torsion, almost torsionfree objects, then Proposition \ref{prop:finitely-cogenerated} would imply that \(\mathbf t\) is finitely generated, hence compact. This is a contradiction. The statement (3) is obtained by duality.
\end{proof}

\subsection{The lattices of torsion and torsionfree classes}

Let \(\torf(\A)\) be the set of torsionfree classes ordered by inclusion and let \(\tors(\A)\) be the set of torsion classes ordered by inclusion.

\begin{theorem}[{\cite[Theorem 3.1]{DIRRT}}]\label{thm:torsionfree-lattice}
The poset \((\torf(\A),\subseteq)\) is a complete lattice. Its meet is given by intersections and its join is given by taking the smallest torsionfree class containing the union. Moreover, this lattice is bi-algebraic and completely semidistributive. Dually, \(\tors(\A)\) is a complete lattice.
\end{theorem}

The following elementary observation allows an immediate dualization of all the lattice theoretical results about torsionfree classes.

\begin{lemma}
The assignment:
\[
	\torf(\A) \to \tors(\A), \quad \mathbf{f} \mapsto {}^{\perp_0}\mathbf{f}
\]
is a duality (i.e. an anti-isomorphism) between the lattice of torsionfree classes and the lattice of torsion classes. 
Moreover the assignment:
\[
	\torf(\A) \to \tors(\A^{\mathrm{op}}), \quad \mathbf{f} \mapsto \mathbf{f}
\]
is a lattice isomorphism.
\end{lemma}

\begin{definition}
An element \( x \) of a complete lattice is said to be \emph{compact} if whenever \( x \le \bigvee_I y_i \) for some collection \( \{ y_i \}_{i \in I} \) of elements of \( L \) there exists a finite subset \( I_0 \subset I \) such that \( x \le \bigvee_{I_0} y_i \). 

Dually, an element \( x \) is \emph{cocompact} if whenever \( x \ge \bigwedge_I y_i \)  for some collection \( \{ y_i \}_{i \in I} \) of elements of \( L \) there exists a finite subset \( I_0 \subset I \) such that \( x \ge \bigwedge_{I_0} y_i \). 
\end{definition}

\begin{remark}
Notice that in a bialgebraic lattice the inequalities in the definition of compact and cocompact elements could be replaced by equalities.

Moreover, by the duality between torsion and torsionfree classes, a torsion class is cocompact if and only if the corresponding torsionfree class is compact. 
\end{remark}

\begin{proposition}[{\cite[Proposition 3.2]{DIRRT}}]\label{prop:compact-finitely-generated}
A torsionfree class is compact in \(\torf(\A)\) if and only if it is finitely cogenerated. Dually, a torsion class is compact in \(\tors(\A)\) if and only if it is finitely generated.
\end{proposition}

\begin{proposition}[{\cite[Theorem 1.0.2, Theorem 2.3.2]{BCZ}}]\label{prop:brick-labels-covers}
Let \((\mathbf t,\mathbf f)\) be a torsion pair.
\begin{enumerate}[label=\textup{(\arabic*)},leftmargin=2.5em]
\item Upper covers of \(\mathbf t\) in \(\tors(\A)\) are in bijection with (isomorphism classes of) torsionfree, almost torsion objects for \((\mathbf t,\mathbf f)\).
\item Lower covers of \(\mathbf t\) in \(\tors(\A)\) are in bijection with (isomorphism classes of) torsion, almost torsionfree objects for \((\mathbf t,\mathbf f)\).
\end{enumerate}
If \(\mathbf t\lessdot \mathbf u\) is a cover relation, the corresponding object is called the brick label of the cover. It is torsionfree, almost torsion for \(\mathbf t\) and torsion, almost torsionfree for \(\mathbf u\).
\end{proposition}

We will use the following notation for labels.

\begin{definition}
For a torsion pair \((\mathbf t,\mathbf f)\), let
\[
\Up(\mathbf t)=\{\text{torsionfree, almost torsion objects for }(\mathbf t,\mathbf f)\}/\cong,
\]
and let
\[
\Low(\mathbf t)=\{\text{torsion, almost torsionfree objects for }(\mathbf t,\mathbf f)\}/\cong.
\]
Thus \(\Up(\mathbf t)\) is the set of upper labels of \(\mathbf t\), while \(\Low(\mathbf t)\) is the set of lower labels of \(\mathbf t\).

We will sometimes abuse notation, for instance writing $ R \in \Up(\mathbf{t}) $ for an object \( R \) meaning that \( R \) is torsionfree, almost torsion for \( \tpair{t}{f} \).
\end{definition}

We now recall a fundamental result and notion from the work of Asai which leads to some sort of mutation process in the context of abelian length category:

\begin{definition}
Let \((\mathbf t,\mathbf f)\) and \((\mathbf u,\mathbf v)\) be torsion pairs with \( \mathbf{t} \subsetneq \mathbf{u} \), we say that they form a \emph{wide interval} if \( \mathbf{f} \cap \mathbf{u} \) is a wide subcategory.
\end{definition}

\begin{remark}
Notice that if \( \mathbf{u} \) covers \( \mathbf{t} \), then the corresponding interval is wide: in fact it has the form \( \filt(B) \) for the brick \( B \) labeling the cover ( see for instance \cite[Theorem 3.3]{DIRRT} ).
\end{remark}

We will use the following notation:

\begin{definition}
For two full subcategories \( \mathcal{L} \) and \( \mathcal{R} \) of \( \A \) we fix
\[
\mathcal{L} * \mathcal{R} := \left\{ M \in \A \,|\,  0 \to L \to M \to R \to 0, \text{ with }L \in \mathcal{L}\text{ and }R \in \mathcal{R} \right\}
\]
\end{definition}

\begin{proposition}[{\cite[Theorem 4.2, Proposition 6.3]{AsaiPfeifer}}]\label{prop:cover-intersection}
Let \((\mathbf t,\mathbf f)\) and \((\mathbf u,\mathbf v)\) be torsion pairs forming a wide interval and let \( \mathcal{W} = \mathbf{u} \cap \mathbf{f} \). Then
\[
\mathbf{t} = \mathbf{u} \cap {}^{\perp_0} \mathcal{W},
\qquad
\mathbf f=\mathcal{W}*\mathbf v.
\]
Dually, 
\[
\mathbf {u} = \mathbf{t} * \mathcal{W}, \qquad \mathbf{f} \cap \mathcal{W}^{\perp_0} = \mathbf{v}
\]
Moreover, the simple objects of \( \mathcal{W} \) are upper labels for \( \mathbf{t} \) and lower labels for \( \mathbf{u} \).
\end{proposition}

An immediate corollary is the characterisation of atoms and co-atoms in the lattice of torsion/torsionfree classes:

\begin{corollary}[{\cite[Proposition 3.5]{AsaiPfeifer}}]
The torsion classes covering the zero torsion class are in bijection with the set of isoclasses of simple objects of \( \mathcal{A} \). 

In particular all such classes have the form \(\filt(S) \) for some simple object $ S $.
\end{corollary}

We finally recall the standard finiteness criterion from \cite{DIJ}, extended to abelian length categories, for instance by Enomoto:

\begin{theorem}[{\cite[Theorem 5.5]{EnomotoMonobrick}}]\label{thm:brick-finite-torsion-finite}
An abelian length category is brick-finite (that is, it contains a finite number of isomorphism classes of bricks) if and only if the set of torsion pairs is finite.
\end{theorem}

\section{Exchange of labels across a wide interval}

 Throughout the section we fix two torsion pairs
\[
(\mathbf t,\mathbf f),\qquad (\mathbf u,\mathbf v),
\]
with
\[
\mathbf t\subsetneq \mathbf u,
\]
being a wide interval 
and we denote by \(\mathcal{W} = \mathbf{u} \cap \mathbf{f} \) the corresponding wide subcategory. Thus as we have recalled the simple objects of \(\mathcal{W}\) are torsionfree, almost torsion for \((\mathbf t,\mathbf f)\) and torsion, almost torsionfree for \((\mathbf u,\mathbf v)\).

It is a natural question to ask whether also the other labels of the lower torsion class are in any relation with labels of the upper torsion class. 

The results contained in this section are probably well-known to experts, at least in the context of finite-dimensional algebras. The collection of upper and lower labels of a given torsion class forms what is called a \emph{semibrick pair} in \cite{BH}, where mutation of such collections is also described in the derived category, in the context of  \( \tau-\)tilting finite algebras.

More recently, the same kind of explicit mutation procedure which we will use has been described in \cite[Theorem 2.5]{MutSimple} with an in depth discussion of mutability criteria in terms of existence of certain minimal left approximations.

Without referring to the derived category, we show that the necessary approximations exist  if we assume that  the torsion pair \( \tpair{t}{f} \) is widely generated and \( \tpair{u}{v} \) is widely cogenerated.

We start in the next subsection with a way to relate upper labels of \( \mathbf{u} \) with labels of \( \mathbf{t} \).

\subsection{A correspondence for upper labels}

Recall the following definition:

\begin{definition}
Let \(\mathcal C\) be a full subcategory of \(\A\). A morphism \(e:X\to C\), with \(C\in\mathcal C\), is a \(\mathcal C\)-envelope if every morphism \(X\to C'\), with \(C'\in\mathcal C\), factors through \(e\), and every endomorphism \(h:C\to C\) with \(he=e\) is an automorphism.

Dually, one defines \( \mathcal C\)-covers.
\end{definition}

We will use several times the following elementary observation:

\begin{lemma}
\label{lem:approxim-lem}
Let $ \B $ be a semibrick and let $ \mathcal{C} = \filt(\B) $. Consider the short exact sequence
\[
0 \to X \to C \to Y \to 0
\]
with $ C \in \mathcal{C} $. Then if $ \Hom(Y, \B) = \Ext^1(Y, \B) = 0 $, the map $ X \to C $ is a $ \mathcal{C}-$envelope.

Moreover in a short exact sequence of the form:
\[
0 \to X \to Y \to C \to 0
\]
with $ C \in \mathcal{C} $, if $ \Hom(X, \B) = 0 $ the map $ Y \to C $ is a $ \mathcal{C}-$envelope.
\end{lemma}

\begin{proof}
As every object $ N \in \mathcal{C} $ admits a finite filtration with factors isomorphic to bricks in $ \B $, we deduce that for all such objects  $ \Hom(Y, N) = \Ext^1(Y, N) = 0 $. Applying the functor $ \Hom(-, N) $ to the short exact sequence in the statement we get an isomorphism $ \Hom(C, N) \simeq \Hom(X, N) $.  

This shows that every map $ X \to N $ factors through the given map $ f :  X \to C $. Moreover, if for some map $ e : C \to C $ we have $ e \circ f = f $, that is $ (1_C - e) \circ f = 0 $, but then as the map $ \Hom(C, C) \to \Hom(X, C) $ is an isomorphism, and $ 0 \circ f = 0 $ we must have $ e = 1_C $, in particular it is an automorphism.

The same argument applies for the second sequence, after computing the long exact sequence given by the functor $ \Hom(-, N) $.
\end{proof}

The exchange process is based on the following Proposition whose proof suggests a label mutation algorithm:

\begin{proposition}
\label{prop:exchange-algo}
Assume that \(\mathbf t\) is widely generated. Let \(V\)  be torsionfree, almost torsion for \(\tpair{u}{v}\), and suppose that there is an exact sequence
\[
\begin{tikzcd}
0\ar[r] & K \ar[r] & V\ar[r] & M \ar[r] & 0
\end{tikzcd}
\]
with \(M\in\mathcal{W}\) and with \(K\) a \(\mathbf v\)-simple. Then there exists a \( \mathcal{W}-\)envelope \( e : V \to W_V \) and one of the following occurs
\begin{enumerate}
\item The morphism \(e \) is an epimorphism and \(\ker e \) is a torsionfree, almost torsion object for \( \tpair{t}{f} \).
\item The morphism \(e \) is a monomorphism and \( \coker e \) is a torsion, almost torsionfree object for \( \tpair{t}{f} \).
\end{enumerate}
\end{proposition}

\begin{proof}
We will do an induction on the length of the kernel \( K \) and divide the proof in three steps. The first two steps are independent of the length of \( K \). Let \( \B \) be the semibrick corresponding to \( \mathcal{W} \).

\noindent\textbf{Step 1}: If \( \Hom(K, \B) = 0 \) then the map \( V \to M \) is already a \( \mathcal{W}-\) envelope by Lemma \ref{lem:approxim-lem}. It remains to show that \( K \) is torsionfree, almost torsion with respect to \( \tpair{t}{f} \). 

First, we show that it is \( \mathbf{f}-\)simple. It is non-zero and contained in $ \mathbf{v} \subset \mathbf{f} $. So consider a proper quotient \( K \to Q \). As \( K \) is \( \mathbf{v}-\)simple, \( Q \in \mathbf{u} \). However, \( K \in {}^{\perp_0} \mathcal{W} \), which is a torsion class. So \( Q \in \mathbf{u} \cap {}^{\perp_0} \mathcal{W} \) which is  \( \mathbf{t} \) by Proposition \ref{prop:cover-intersection}.

Now we need to show maximality among \( \mathbf{f}-\)simples. So let $ K \to K' $ be a non-zero map to a relative simple, this map is necessarily a monomorphism. Assume it is proper. Form the pushout:
\[
\begin{tikzcd}
0 \ar[r] & K \ar[r] \ar[d, hook] & V \ar[r] \ar[d, hook] & M \ar[r] \ar[d, equals] & 0 \\
0 \ar[r] & K' \ar[d, two heads] \ar[r] & P \ar[d, two heads] \ar[r] & M \ar[r] & 0 \\
& C \ar[r, equals] & C
\end{tikzcd}
\]
Then $ P \in \mathbf{f} $, being an extension of $ M \in \mathcal{W} \subset \mathbf{f} $ and $ K' $. Moreover \( C \in \mathbf{t} \) being a proper quotient of \( K' \).

 Notice that the short exact sequence \( 0 \to V \to P \to C \to 0 \) is non-split, otherwise \( C = 0 \) as \( P \in \mathbf{f} \), thus being \( V \) torsionfree, almost torsion \( P \in \mathbf{u} \).

In conclusion \( P \in \mathbf{f} \cap \mathbf{u} = \mathcal{W} \). But then the map \( K \to  V \to P \) is the zero map, as \( \Hom(K, \mathcal{W}) = 0 \), which is a contradiction as it is the composite of two non-zero monomorphisms. This shows that the monomorphism \( K \to K' \) can't be proper, so \( K \) is maximal among \( \mathbf{f}-\)simples.

\noindent \textbf{Observation}: If \( f : K \to B \) is non-zero, for some \( B \in \mathcal{B} \) then it is either a (proper) monomorphism or a (proper) epimorphism. 

It clearly can't be an isomorphism, as \( K \in \mathbf{v} \) and \( B \in \mathbf{u} \). So assume it has a non-zero kernel, then the image of \( f \) is a proper quotient of \( K \), thus it lies in \( \mathbf{u} \) and is non-zero. But every proper subobject of \( B \) lies in \( \mathbf{v} \) as \( B \) is \( \mathbf{u}-\)simple. Thus \( f \) must be a proper epimorphism. 

\noindent \textbf{Step 2}: Assume we have a proper monomorphism \( K \to B \) for some \( B \in \B \). Then consider once again the pushout diagram
\[
\begin{tikzcd}
0 \ar[r] & K \ar[r] \ar[d, hook] & V \ar[r] \ar[d, hook] & M \ar[r] \ar[d, equals] & 0 \\
0 \ar[r] & B \ar[d, two heads] \ar[r] & P \ar[d, two heads] \ar[r] & M \ar[r] & 0 \\
& Q \ar[r, equals] & Q
\end{tikzcd}
\]
Notice that \( P \in \mathcal{W} \). I claim that \( Q \) is \( \mathbf{t}-\)simple. It is non-zero and contained in \( \mathbf{t} \) being a proper quotient of \( B \). Now it is enough to show that it has no proper non-zero subobject in \( \mathbf{t} \). Assume \( T \to Q \) is a non-zero subobject. Then, pulling back along \( P \to Q \) we obtain:
\[
\begin{tikzcd}
0 \ar[r] & V \ar[r] \ar[d, equals] & R \ar[r] \ar[d, hook] & T \ar[r] \ar[d, hook] & 0 \\
0 \ar[r] & V  \ar[r] & P \ar[d, two heads] \ar[r] & Q \ar[r] \ar[d, two heads] & 0 \\
& & C \ar[r, equals] & C
\end{tikzcd}
\]
Notice that \( R \in \mathbf{f} \) as it is a subobject of  \(  P \in \mathcal{W} \). Moreover, as the sequence \( 0 \to V \to R \to T \to 0 \) is non-split, \( R \in \mathbf{u} \) as \( V \) is torsionfree, almost torsion for \( \tpair{u}{v} \) and \( T \in \mathbf{t} \subset \mathbf{u} \). Thus \( R \in \mathcal{W} \), but as this is a wide subcategory, so is the cokernel \( C \). But \( C \) should be in the torsion class \( \mathbf{t} \) as it is (isomorphic to) a quotient of \( Q \). Thus \( C = 0 \) and \( T \) is not a proper subobject. 

Now \( Q \) may not be torsion, almost torsionfree, however, as we assumed \( \mathbf{t} \) to be widely generated, there exists a torsion, almost torsionfree object \( \widetilde{Q} \) that has \( Q \) as a quotient. 

Construct the pullback
\[
\begin{tikzcd}
& & L \ar[r, equals] \ar[d, hook] & L \ar[d, hook] \\
0 \ar[r] & V  \ar[d, equals] \ar[r] & M' \ar[d, two heads] \ar[r] & \widetilde{Q} \ar[r] \ar[d, two heads] & 0 \\
0 \ar[r] & V  \ar[r] & P \ar[r] & Q \ar[r] & 0
\end{tikzcd}
\]
\( L \) is a possibly zero, proper subobject of a torsion, almost torsionfree, thus it is contained in \( \mathbf{f} \). This means that \( M' \) is in \( \mathbf{f} \) as an extension of torsionfree objects, moreover \( M' \in \mathbf{u} \) as part of a non-split short exact sequence with \( \widetilde{Q} \in \mathbf{u} \) and \( V \) torsionfree, almost  torsion as above. This shows \( M' \in \mathcal{W} \) and we obtain that \( V \to M' \) is the required envelope using Lemma \ref{lem:approxim-lem} after noticing that the required orthogonality conditions hold for the torsion, almost torsionfree object \( \widetilde{Q} \) and the simple objects of \( \mathcal{W} \) which are torsionfree, almost torsion for the same torsion pair. 

\noindent \textbf{Step 3}: Here we use the length of \( K \). If it is minimal (that is 1, as \( K \ne 0 \) ) then \( K \) is simple and we can't have a proper epimorphism \( K \to B \) for any \( B \in \B \). Thus this completes the base case of induction. So assume the algorithm works for lengths up to \( n \) and assume we have a \( K \) with a proper epimorphism onto some \( B \). We will construct a new short exact sequence \( 0 \to K' \to V \to M' \to 0 \) with \( K' \) being \( \mathbf{v}-\)simple of strictly lower length and \(M' \in \mathcal{W} \), so that we can conclude by induction hypothesis.  

This is once again obtained by taking a pushout of the epimorphism \( K \to B \):
\[
\begin{tikzcd}
 & K' \ar[r, equals] \ar[d, hook] & K' \ar[d, hook] \\
0 \ar[r] & K \ar[d, two heads] \ar[r] & V \ar[d, two heads] \ar[r] & M \ar[r] \ar[d, equals] & 0 \\
0 \ar[r] & B  \ar[r] & M' \ar[r] & M \ar[r] & 0
\end{tikzcd}
\] 
Notice that \( M' \in \mathcal{W} \) as required and \( \ell(K') < \ell(K) \). So it is enough to show that \( K' \) is \( \mathbf{v}-\)simple. But this is immediate: take a proper quotient \( S \) of \( K' \), once again this induces a short exact sequence \( 0 \to S \to S' \to B \to 0 \), where \( S' \in \mathbf{u} \) being a proper quotient of the \( \mathbf{v}-\)simple \( K \). Thus \( S \) is in \( \mathbf{u} \) as \( B \) is torsion, almost torsionfree for \( \tpair{u}{v} \). 
\end{proof}

\begin{corollary}
\label{cor:up-exchange}
Assume that \(\mathbf t\) is widely generated. Let \(V\)  be torsionfree, almost torsion for \(\tpair{u}{v}\)
Then there exists a \( \mathcal{W}-\)envelope \( e : V \to W_V \) and one of the following occurs
\begin{enumerate}
\item The morphism \(e \) is an epimorphism and \(\ker e \) is a torsionfree, almost torsion object for \( \tpair{t}{f} \).
\item The morphism \(e \) is a monomorphism and \( \coker e \) is a torsion, almost torsionfree object for \( \tpair{t}{f} \).
\end{enumerate}
\end{corollary}

\begin{proof}
Apply Proposition \ref{prop:exchange-algo} to the trivial short exact sequence \[ 0 \to V \to V \to 0 \to 0. \] Of course \( V \) is \( \mathbf{v}-\)simple and \( 0 \in \mathcal{W} \).
\end{proof}

\begin{remark}
Notice that in the corollary above, if \( \Hom(V, \mathcal{W}) = 0 \) then the envelope is the zero map and \( V \) itself is a torsionfree, almost torsion object for \( \tpair{t}{f} \).
\end{remark}

\subsection{Exchange for lower labels}

We now show how to associate lower labels of \( \mathbf{u} \) to lower labels of \( \mathbf{t} \).



\begin{proposition}
\label{prop:lower-exchange}
Assume that \(\mathbf t\) is widely generated. Let \(U\)  be torsion, almost torsionfree for \(\tpair{u}{v}\). Then if \( U \not \in \mathcal{W} \), there exists a short exact sequence
\[
0 \to W_T \to T \to U \to 0 
\]
where the map \( W_T \to T \) is a \( \mathcal{W}-\)cover and \( T \) is torsion, almost torsionfree for \( \tpair{t}{f} \). Moreover, if \( T' \) is a torsion, almost torsionfree object for \( \tpair{t}{f} \) not isomorphic to \( T \), we have \( \Hom(T', U) = 0 \).
\end{proposition}

\begin{proof}
First, notice that all the torsion, almost torsionfree objects of \( \tpair{u}{v} \) not contained in \( \mathcal{W} \) are contained in \( {}^{\perp_0}\mathcal{W} \cap \mathbf{u} = \mathbf{t} \) ( see Proposition \ref{prop:cover-intersection}) and are in particular, \( \mathbf{t}-\)simple, as all their proper subobjects lie in \( \mathbf{v} \subset \mathbf{f} \). 

Now, if such an object \( U \) is already a torsion, almost torsionfree object for \( \tpair{t}{f} \), then we are done, as the trivial short exact sequence \( 0 \to 0 \to U \to U \to 0 \) has the claimed properties and Hom-orthogonality to non-isomorphic torsion, almost torsionfree object is immediate by Proposition \ref{prop:maximal-f-simple}. 

If it is not, then as \( \mathbf{t} \) is widely generated, we have some torsion, almost torsionfree object \( T \) that has \( U \) as a proper quotient. Consider the corresponding short exact sequence
\[
0 \to K \to T \to U \to 0
\]
As \( U \) is torsion, almost torsionfree for \( \tpair{u}{v} \), the kernel lies in \( \mathbf{u} \), but as \( T \) is in particular \( \mathbf{t}-\)simple \( K \in \mathbf{f} \). So \( K \in \mathcal{W} \) as required.
As \( \Hom(\mathcal{W}, U) = 0 \), the map \( K \to T \) is a \( \mathcal{W}-\)cover, by the dual of Lemma \ref{lem:approxim-lem}.

Now, let \( T' \) be a torsion, almost torsionfree object, not-isomorphic to \( T \). Then, applying the functor \( \Hom(T', -) \) to the approximation sequence we get the exact sequence
\[
\Hom(T', T) \to \Hom(T', U) \to \Ext^1(T', K)
\]
but \( \Hom( T', T) = 0 \) as these are non-isomorphic lower labels, moreover \( \Ext^1(T', B) = 0 \) for all the simple objects of the wide subcategory \( \mathcal{W} \) by Lemma \ref {lem:ext-orthogonality} so that also \( \Ext^1(T', K) = 0 \) as \( K \in \mathcal{W} \) is filtered by the simple \(\mathcal{W}-\)objects.
Thus \( \Hom(T', U) = 0 \).
\end{proof}

\subsection{The bijection of labels}

We write down explicitly the dual results before giving the full exchange theorem:

\begin{corollary}
\label{cor:up-exchange-dual}
Assume that \(\mathbf v\) is widely cogenerated. Let \(T\)  be torsion, almost torsionfree for \(\tpair{t}{f}\)
Then there exists a \( \mathcal{W}-\)cover \( c : W_T \to T \) and one of the following occurs
\begin{enumerate}
\item The morphism \(c \) is an epimorphism and \(\ker c \) is a torsionfree, almost torsion object for \( \tpair{u}{v} \).
\item The morphism \(c \) is a monomorphism and \( \coker c \) is a torsion, almost torsionfree object for \( \tpair{u}{v} \).
\end{enumerate}
\end{corollary}

\begin{proposition}
\label{prop:upper-exchange-dual}
Assume that \(\mathbf v\) is widely cogenerated. Let \(F\)  be torsionfree, almost torsion for \(\tpair{t}{f}\). Then if \( F \not \in \mathcal{W} \), there exists a short exact sequence
\[
0 \to F \to V \to W_V \to 0 
\]
where the map \( V \to W_V \) is a \( \mathcal{W}-\)envelope and \( V \) is torsionfree, almost torsion for \( \tpair{u}{v} \). Moreover, if \( V' \) is a torsionfree, almost torsion object for \( \tpair{u}{v} \) not isomorphic to \( V \), we have \( \Hom(F, V') = 0 \).
\end{proposition}


\begin{theorem}
\label{thm:full-mutation-labels}
Let \( \tpair{t}{f} \) and \( \tpair{u}{v} \) be torsion pairs with \( \mathbf{t} \subsetneq \mathbf{u} \) a wide interval and such that the lower pair is widely generated, while the upper pair is widely cogenerated.
Then there is a bijection
\[
\Phi : \Up(\mathbf u)\coprod \Low(\mathbf u) \to \Up(\mathbf t)\coprod \Low(\mathbf t).
\]
defined as follows:
\begin{enumerate}
\item for an upper label $ [V] $ of $ \mathbf{u} $, take a $ \mathcal{W}-$envelope $ e_V : V \to W_V $. This is either a monomorphism or an epimorphism, if mono, $ \Phi([V]) := [\coker e_V] $ if epi $ \Phi([V]) := [\ker e_V] $.
\item for a lower label $ [U] $, if $ U \in \mathcal{W} $, then $ \Phi([U]) := [U] $ otherwise we fix $ \Phi([U]) := [T] $ where $ T $ is the unique torsion, almost torsionfree object with \( \Hom(T, U) \ne 0 \).
\end{enumerate}
\end{theorem}

\begin{proof}
First we notice that the map is well defined: the associated objects exist and are upper and lower labels of \( \tpair{t}{f} \) by Corollary \ref{cor:up-exchange} and  Proposition \ref{prop:lower-exchange}, moreover isomorphic objects have isomorphic envelopes and via the factorisation property, isomorphic kernels/cokernels.

To show it is a bijection, we consider the following map and show it is its inverse.
Let \( \Psi :  \Up(\mathbf t)\coprod \Low(\mathbf t) \to  \Up(\mathbf u)\coprod \Low(\mathbf u) \) be defined as follows:
\begin{enumerate}
\item  for a lower label $ [T] $, take  a \(\mathcal{W}-\)cover \( c_T : W_T \to T \).  This is either a monomorphism or an epimorphism, if mono, $ \Psi([T]) := [\coker c_T] $ if epi $ \Psi([T]) := [\ker c_T] $.
\item for an upper label \( [F] \), if \( F \in \mathcal{W} \) then \( \Psi([F]) = [F] \), otherwise we fix \( \Psi([F]) = [V] \) where \( V \) is the unique (up to isomorphism) torsionfree, almost torsion object for \( \tpair{u}{v} \) admitting a non-zero map from \( F \).
\end{enumerate}

Once again, this map is well defined by Corollary \ref{cor:up-exchange-dual} and  Proposition \ref{prop:upper-exchange-dual}. 

Let's compute \( \Psi(\Phi([V])) \) for an upper label of \( \tpair{u}{v} \). If the envelope is a monomorphism, then we get a lower label \(  [\coker e_V] \) of \( \tpair{t}{f} \). But observe that the short exact sequence
\[
0 \to V \to W_V \to \coker e_V \to 0
\]
also gives a \( \mathcal{W}-\)cover of \( \coker e_V \) by the dual of Lemma \ref{lem:approxim-lem}. So \( [\coker e_V] \) is a lower label with epimorphic cover and \( \Psi( [\coker e_V] ) = [ V] \).

On the other hand, if the envelope is an epimorphism, we have the short exact sequence
\[
0 \to \ker e_V \to V \to W_V \to 0
\]
and \( [ \ker e_V ] \) is an upper label for \( \tpair{t}{f} \). Notice that \( \ker e_V \not \in \W \) as \( \Hom(\mathcal{W}, V) = 0 \), moreover \( V \) is the only torsionfree, almost torsion object up to isomorphism with \( \Hom( \ker e_V, V) \ne 0 \). This is clear, applying to the short exact sequence the functor \( \Hom(-, V') \) and using the usual Hom and Ext-orthogonality properties. Thus \( \Psi([\ker e_V ]) = [V] \). 

Let's now compute  \( \Psi(\Phi([U])) \) for a lower label of \( \tpair{u}{v} \). If \( [U] \in \mathcal{W} \) then we are done. So assume \( U \not \in \mathcal{W} \) and let \( [T] \) be the associated lower label of \( \tpair{t}{f} \). In every short exact sequence
\[
0 \to W_T \to T \to U \to 0
\] 
we now know that the map \( W_T \to T \) is a \( \mathcal{W}-\)cover as in the proof of Proposition \ref{prop:lower-exchange}. But then \( \Psi([T]) = [U] \). 

This shows that \( \Psi \circ \Phi = \mathrm{Id} \). One checks, by dual computations that also \( \Phi \circ \Psi = \mathrm{Id} \).
\end{proof}


\begin{thebibliography}{99}
\bibitem[AsaiPfeifer]{AsaiPfeifer} S. Asai and C. Pfeifer; \emph{Wide subcategories and lattices of torsion classes}, Algebr. Represent. Theory 25 (2022), no. 6, 1611--1629.
\bibitem[BCZ]{BCZ} E. Barnard, A. Carroll and S. Zhu; \emph{Minimal inclusions of torsion classes}, Algebr. Comb. 2 (2019), no. 5, 879--901.
\bibitem[BH]{BH} E. Barnard, E. J. Hanson; \emph{Pairwise compatibility for 2-simple Minded collections II: Preprojective algebras and Semibrick Pairs of Full Rank}, Ann. Comb. 26, 803-855 (2022).
\bibitem[DIJ]{DIJ} L. Demonet, O. Iyama and G. Jasso; \emph{\(\tau\)-tilting finite algebras, bricks and \(g\)-vectors}, Int. Math. Res. Not. IMRN 2019, no. 3, 852--892.
\bibitem[DIRRT]{DIRRT} L. Demonet, O. Iyama, N. Reading, I. Reiten and H. Thomas; \emph{Lattice theory of torsion classes: beyond \(\tau\)-tilting theory}, Trans. Amer. Math. Soc. Ser. B 10 (2023), 542--612.
\bibitem[EnomotoMonobrick]{EnomotoMonobrick} H. Enomoto; \emph{Monobrick, a uniform approach to torsion-free classes and wide subcategories}, Adv. Math. 393 (2021), Paper No. 108113.
\bibitem[MutSimple]{MutSimple} S. Asai; \emph{Mutations of simple-minded collections revisited}; preprint, arXiv:2606.15537, (2026)
\bibitem[RingelBrickChain]{RingelBrickChain} C. M. Ringel; \emph{Brick chain filtrations}, preprint, arXiv:2411.18427v2 (2024).
\bibitem[RingelSpecies]{RingelSpecies} C. M. Ringel; \emph{Representations of \(K\)-species and bimodules}, J. Algebra 41 (1976), no. 2, 269--302.
\end{thebibliography}
\end{document}
